\section{ROUND-4 OBJECTIVE}

I continue along path \textbf{(C)}: strengthen the Round-3 obstruction for the minimally corrected statement.

Round 2 already identified the fatal ambiguity in the naive formulation: without distinctness of the moduli, the question is trivial, since
\[
\{x \equiv 0 \pmod 3,\ x \equiv 1 \pmod 3,\ x \equiv 2 \pmod 3\}
\]
is an odd covering system. Thus the genuine problem remains the standard one:
\begin{quote}
Does there exist a \emph{distinct} covering system all of whose moduli are odd integers greater than \(1\)?
\end{quote}

Round 3 proved that any such system must contain a modulus at least \(127\). In this round I strengthen that finite obstruction substantially:

\begin{quote}
\textbf{Main new theorem.} Any distinct odd covering system must contain a modulus at least \(1001\).
\end{quote}

This is still far from a full solution, but it is a genuine theorem strictly stronger than the Round-3 lower bound.

\section{Round-3 FOUNDATION USED}

I use the following Round-3 results as fixed input.

\begin{enumerate}
\item \textbf{Corrected formulation.} The nontrivial problem is the \emph{distinct} odd-modulus problem; without distinctness the answer is trivially ``yes.'' 
\item \textbf{Round-3 prime-power assignment obstruction (Lemma 3.1).} If a finite set \(M\) of odd moduli is assigned to prime divisors \(\pi(m)\mid m\) in such a way that, for every prime \(p\),
\[
\lambda_p:=\sum_{\pi(m)=p} p^{-v_p(m)} < 1,
\]
then for \emph{every} choice of residue classes \(x\equiv a_m \pmod m\) with \(m\in M\), there exists an integer \(x\) avoiding all of them. In particular, the chosen congruences do not cover \(\mathbb Z\).
\item \textbf{Round-3 lower bound.} Any distinct odd covering system must contain a modulus at least \(127\).
\end{enumerate}

I do not re-prove these items here. The new work below uses item 2 and strictly strengthens item 3.

\section{NEW INSIGHT / TOOL (ROUND-4)}

The new idea is that the Round-3 prime-assignment lemma can be pushed much further by using the following assignment rule.

\begin{itemize}
\item Default rule: assign each odd modulus \(m\le 999\) to its \emph{largest prime divisor}.
\item Then make only \(23\) explicit exceptions, all involving moduli whose prime divisors are at most \(19\).
\end{itemize}

This assignment was \emph{found} by an exact mixed-integer feasibility search, but the proof below does not depend on the search. Once the assignment is written down, the verification is completely explicit: one checks that every local load \(\lambda_p\) is strictly less than \(1\), and then Round-3 Lemma 3.1 applies.

So the genuinely new content of this round is an explicit certificate showing that the Round-3 obstruction already rules out \emph{all} odd moduli up to \(999\).

\section{ATTACK PLAN (ROUND-4)}

After Round 3, the exact remaining gap was the following: the best unconditional finite obstruction only showed that a hypothetical distinct odd covering system must contain some modulus at least \(127\).

To improve this, I will:

\begin{enumerate}
\item use Round-3 Lemma 3.1 as the main engine;
\item define a prime-divisor assignment \(\pi(m)\) for every odd \(m\) with \(3\le m\le 999\);
\item verify that every local load \(\lambda_p\) is \(<1\);
\item conclude that no cover can use only odd moduli at most \(999\), hence any distinct odd covering system must contain a modulus at least \(1001\).
\end{enumerate}

This overcomes the Round-3 limitation because the new assignment is much more efficient than the previous explicit partitions: almost every modulus is handled by its largest prime divisor, and only a small exceptional family has to be rerouted.

\section{WORK (ROUND-4)}

\subsection*{Theorem 4.1}
Let
\[
M_{999}:=\{m\in \mathbb N : 3\le m\le 999,\ m\ \text{odd}\}.
\]
Then for \emph{every} choice of residue classes \(a_m \pmod m\) with \(m\in M_{999}\), there exists an integer \(x\) such that
\[
x \not\equiv a_m \pmod m \qquad (m\in M_{999}).
\]
Consequently, no distinct odd covering system can have all of its moduli at most \(999\). In particular, any distinct odd covering system, if one exists, must contain a modulus at least \(1001\).

\medskip
\noindent\emph{Proof.}
By Round-3 Lemma 3.1, it is enough to define a map
\[
\pi:M_{999}\to \{\text{odd primes}\}
\]
such that \(\pi(m)\mid m\) for every \(m\in M_{999}\) and
\[
\lambda_p:=\sum_{\pi(m)=p} p^{-v_p(m)} < 1
\qquad\text{for every odd prime }p.
\]

For an odd integer \(m\), let \(P(m)\) denote its largest prime divisor. Define \(\pi(m)\) by the default rule
\[
\pi(m)=P(m),
\]
except for the following \(23\) moduli:
\[
\begin{aligned}
&45,63,135,189,297,351,405,567,675,891,945 \mapsto 3,\\
&165,175,275,325,525,825,875,975 \mapsto 5,\\
&539,637,833 \mapsto 7,\\
&429 \mapsto 11.
\end{aligned}
\]
Each displayed reassignment is to a genuine prime divisor of the modulus, so \(\pi(m)\mid m\) always holds.

We now verify that every local load is \(<1\).

\subsubsection*{Large primes \(p\ge 23\)}
If \(p\ge 23\), then no exception is assigned to \(p\), so \(\pi(m)=p\) means exactly that \(p=P(m)\), i.e. \(m\) is an odd multiple of \(p\) whose prime factors are all at most \(p\). In particular, every such \(m\) is an odd multiple of \(p\) not exceeding \(999\).

The odd multiples of \(p\) up to \(999\) are
\[
p,\, 3p,\, 5p,\, \dots,\, (2k-1)p
\]
with \((2k-1)p\le 999\). Since \(p\ge 23\),
\[
2k-1 \le \frac{999}{p} \le \frac{999}{23} < 44,
\]
hence \(k\le 22\). Therefore there are at most \(22\) odd multiples of \(p\) in \(M_{999}\).

Every modulus assigned to \(p\) contributes \(p^{-v_p(m)}\le 1/p\). Hence
\[
\lambda_p \le \frac{22}{p} < 1
\qquad (p\ge 23).
\]

\subsubsection*{Small primes \(p\in\{3,5,7,11,13,17,19\}\)}
Under the above assignment rule, the corresponding fibers \(S_p:=\pi^{-1}(p)\) are:

\[
\begin{aligned}
S_3&=\{3,9,27,45,63,81,135,189,243,297,351,405,567,675,729,891,945\},\\
S_5&=\{5,15,25,75,125,165,175,225,275,325,375,525,625,825,875,975\},\\
S_7&=\{7,21,35,49,105,147,245,315,343,441,539,637,735,833\},\\
S_{11}&=\{11,33,55,77,99,121,231,363,385,429,495,605,693,847\},\\
S_{13}&=\{13,39,65,91,117,143,169,195,273,455,507,585,715,819,845\},\\
S_{17}&=\{17,51,85,119,153,187,221,255,289,357,425,459,561,595,663,765,867,935\},\\
S_{19}&=\{19,57,95,133,171,209,247,285,323,361,399,475,513,627,665,741,855,931,969\}.
\end{aligned}
\]

From these explicit sets we read off the \(p\)-adic exponent counts and hence the local loads:

\[
\lambda_3
= \frac{1}{3}+\frac{3}{3^2}+\frac{7}{3^3}+\frac{4}{3^4}+\frac{1}{3^5}+\frac{1}{3^6}
= \frac{715}{729}<1,
\]
\[
\lambda_5
= \frac{3}{5}+\frac{9}{5^2}+\frac{3}{5^3}+\frac{1}{5^4}
= \frac{616}{625}<1,
\]
\[
\lambda_7
= \frac{5}{7}+\frac{8}{7^2}+\frac{1}{7^3}
= \frac{302}{343}<1,
\]
\[
\lambda_{11}
= \frac{10}{11}+\frac{4}{11^2}
= \frac{114}{121}<1,
\]
\[
\lambda_{13}
= \frac{12}{13}+\frac{3}{13^2}
= \frac{159}{169}<1,
\]
\[
\lambda_{17}
= \frac{16}{17}+\frac{2}{17^2}
= \frac{274}{289}<1,
\]
\[
\lambda_{19}
= \frac{18}{19}+\frac{1}{19^2}
= \frac{343}{361}<1.
\]

Thus \(\lambda_p<1\) for every odd prime \(p\).

Round-3 Lemma 3.1 now applies, so for every choice of congruences
\[
x\equiv a_m \pmod m \qquad (m\in M_{999}),
\]
there is an integer \(x\) avoiding all of them. Therefore the full family \(M_{999}\) cannot support a cover.

Finally, if some distinct odd covering system had all of its moduli at most \(999\), then its moduli would form a subset \(M\subseteq M_{999}\). Restricting \(\pi\) from \(M_{999}\) to \(M\) only decreases each local load \(\lambda_p\), so the same argument shows that \(M\) also cannot support a cover. This contradiction proves that no distinct odd covering system can have all moduli at most \(999\). Hence any such system must contain a modulus at least \(1001\). \hfill\(\square\)

\medskip
\noindent\textbf{Remark.}
This strictly strengthens the Round-3 bound
\[
\max m_i \ge 127
\]
to the new unconditional bound
\[
\max m_i \ge 1001.
\]

\section{ADVERSARIAL VERIFICATION}

I check the new argument against the main failure modes.

\begin{enumerate}
\item \textbf{Use of Round-3 results.} The only imported result is Round-3 Lemma 3.1, used exactly as intended: once a prime-divisor assignment with all local loads \(<1\) is produced, noncoverage follows.

\item \textbf{Assignment validity.} Every exceptional reassignment is to an actual prime divisor:
\[
\begin{aligned}
&45,135,405,675,975 \text{ are divisible by }5\text{ and }3;\\
&63,189,567,945 \text{ are divisible by }7\text{ and }3;\\
&297,891 \text{ are divisible by }11\text{ and }3;\\
&351 \text{ is divisible by }13\text{ and }3;\\
&165,825 \text{ are divisible by }11\text{ and }5;\\
&175,525,875 \text{ are divisible by }7\text{ and }5;\\
&275 \text{ is divisible by }11\text{ and }5;\\
&325,975 \text{ are divisible by }13\text{ and }5;\\
&539,637,833 \text{ are divisible by }7;\\
&429 \text{ is divisible by }11.
\end{aligned}
\]

\item \textbf{Large-prime estimate.} For \(p\ge 23\), the proof uses only the crude bound ``at most \(22\) odd multiples of \(p\) up to \(999\)'', together with the trivial inequality \(p^{-v_p(m)}\le 1/p\). This is independent of any hidden structure and gives a strict inequality \(\lambda_p<1\).

\item \textbf{Small-prime loads.} The sets \(S_3,S_5,S_7,S_{11},S_{13},S_{17},S_{19}\) are explicit; the exponent counts used in the displayed formulas are read directly from those lists. Every resulting rational number is strictly below \(1\).

\item \textbf{Subset monotonicity.} The conclusion for arbitrary distinct odd modulus sets with \(\max m_i\le 999\) is legitimate, because restricting the assignment \(\pi\) from \(M_{999}\) to a smaller set can only decrease the local loads. So if the full family cannot cover, then no subfamily can cover either.

\item \textbf{Interaction with Round 3.} There is no conflict: Round 3 gave the weaker lower bound \(127\le \max m_i\), and Theorem 4.1 sharpens it to \(1001\le \max m_i\).
\end{enumerate}

No gap remains in the proof of Theorem 4.1.

\section{FINAL}

\textbf{UNRESOLVED (BUT STRICTLY ADVANCED)}

The corrected odd-covering problem remains open, but this round proves a substantially stronger finite obstruction than Round 3:

\begin{quote}
If a distinct odd covering system exists, then it must contain a modulus at least \(1001\).
\end{quote}

The new ingredient is an explicit prime-divisor assignment certificate on all odd moduli \(3\le m\le 999\), verified using the Round-3 prime-power obstruction lemma.

\section{COMPLETION ESTIMATE}

\textbf{COMPLETION: 58\%}

\section{REFERENCES}

\begin{enumerate}
\item T. F. Bloom, \emph{Erd\H{o}s Problem \#7}, \texttt{https://www.erdosproblems.com/7}, accessed 2026-03-07.
\item C. Bispels, M. Cohen, J. Harrington, J. Lowrance, K. Pontes, L. Schaumann, and T. W. H. Wong, \emph{A further investigation on covering systems with odd moduli}, \emph{Discrete Mathematics} \textbf{349} (2026), no.~7, 115013.
\end{enumerate}
